\cotefile{106}
\scanpage{2}
% slanted
\margin{$\mathrm{Ho}\, C$ pour $C$ une cat. à cofibrations. ---}
\margin{?}

Soit $(C,\mathrm{W}_{\mathrm e},\mathrm{cof})$ une catégorie satisfaisant C1 à C3 de Baues, (pas néc. C4) (p. ex on pourrait avoir $\mathrm{cof}=\unsure{\mathrm{cof}|\mathrm{wcof}}_{\mathrm{we}}$). \inserted{On suppose qu'il y a objet initial $\varnothing$} Alors il est connu que
\[
 W_{\mathrm{Cof}}^{-1}\,\mathrm{Cof}\longrightarrow W^{-1}C
\]
est une équivalence de catégories.

Pour le voir, on va définir un foncteur en sens inverse (qui sera quasi-inverse)
\[
 \Phi\colon W^{-1}C\longrightarrow W_{\mathrm{Cof}}^{-1}\,\mathrm{Cof},\quad
 \text{i.e.}\quad C\longrightarrow W_{\mathrm{Cof}}^{-1}.\mathrm{Cof}
 \quad\text{qui rend les }u\in W\text{ inv.}
\]

Pour tt $x\in C$, utilisant C3 pour $\varnothing\longrightarrow x$, on trouve
\[
 \widetilde x\xrightarrow[\sim]{\alpha_x}x
 \qquad\text{avec }\widetilde x\in\mathrm{Cof},\ \alpha_x\in W .
\]

Soit alors $x\xrightarrow{f}y$ dans $C$,
\diagram{\begin{tikzcd}[ampersand replacement=\&]
\widetilde x \arrow[d,"\sim"',"{\alpha_x}"'{pos=.65}] \arrow[dr,"{f\alpha_x}"] \& \widetilde y \arrow[d,"{\alpha_y}"]\\
x \arrow[r,"f"'] \& y
\end{tikzcd}}
considérons
\diagram{\begin{tikzcd}[ampersand replacement=\&]
\varnothing \arrow[d] \arrow[r] \& \widetilde y \arrow[d,"\sim"' ,"{\alpha_y}"]\\
\widetilde x \arrow[r,"{f\alpha_x}"] \& y
\end{tikzcd}}
appliquant le \lq\lq weak lifting lemma\rq\rq, on trouve
\diagram{\begin{tikzcd}[ampersand replacement=\&,column sep=small]
\& \widetilde y \arrow[r,tail,"i","\sim"'] \arrow[d,"\sim"'," {\alpha_y}"] \& \widetilde y' \arrow[dl,"\sim" ,"{\alpha'_{\deleted{\unsure{y|g}}}}"' ]\\
\widetilde x \arrow[urr,bend left=12,"{\widetilde\varphi}"] \arrow[r,"{f\alpha_x}","f"'] \& y \&
\end{tikzcd}}
un relèvement $\widetilde\varphi$ \deleted{\unsure{vers, pas de}} de $f\alpha_x$ non pas \deleted{$f\alpha_x$} vers $\widetilde y$, mais vers un $\widetilde y'$.

Considérons alors
\[
 [i]^{-1}\widetilde\varphi\colon\widetilde x\longrightarrow\widetilde y
 \qquad\text{dans }\mathrm{Hom}_{W_{\mathrm{Cof}}^{-1}}(\widetilde x,\widetilde y)
\]
Je dis qu'il ne dépend pas du choix du $(i,\alpha'_{\unsure{y|g}},\varphi)$. Pour deux tels choix \deleted{\unsure{on}}
\diagram{\begin{tikzcd}[ampersand replacement=\&]
\& \widetilde y'_1\\
\widetilde y \arrow[ur,tail,"{i_1}"] \arrow[dr,tail,"{i_2}"'] \&\\
\&\widetilde y'_2
\end{tikzcd}}
parallèles : introduire $\widetilde y'_3=\widetilde y'_1\vee_{\widetilde y}\widetilde y'_2$, OPS que $\widetilde y'_1=\widetilde y'_2$, $i_1=i_2$, $\alpha'_1=\alpha'_2$, donc on a deux relèvements $\varphi_1,\varphi_2$ de $f\alpha_x$ vers la $\mathord{=}$ $\widetilde y'$, \deleted{OPS \unsure{$=$|et} que $\widetilde y$} et on est ramené (écrivant $y'$ au lieu de $\widetilde y'$, $\deleted{\unsure{g}}\inserted{g=f\alpha_x}$, $x'$ au lieu de $\widetilde x$, mais $x$ cof.)

\diagram{\begin{tikzcd}[ampersand replacement=\&]
x' \arrow[r,shift left,"\varphi"] \arrow[r,shift right,"\psi"'] \& y' \arrow[r,"\alpha","\sim"'] \& y
\end{tikzcd}}
\deleted{\diagram{\begin{tikzcd}[ampersand replacement=\&]
\& y'\ {\rm cof}\arrow[d,"\alpha","\sim"']\\
x'\ {\rm cof}\arrow[ur,shift left,"\varphi"]\arrow[ur,shift right,"\psi"']\arrow[r,"g"']\&y
\end{tikzcd}}}
prouver que si $x',y'\in\mathrm{Cof}$, et $\alpha\varphi=\alpha\psi\deleted{\unsure{{}=g|\alpha g}}$, alors $[\varphi]=[\psi]$ dans $W_{\mathrm{Cof}}^{-1}\,\mathrm{Cof}$.

\diagram{\begin{tikzcd}[ampersand replacement=\&]
x'\arrow[r]\& y\\
x'\amalg x'\arrow[u,"{\nabla_{x'}}"]\arrow[r,"{(\varphi,\psi)}"']\& y'\arrow[u,"{\in W}"]
\end{tikzcd}}
\diagram{\begin{tikzcd}[ampersand replacement=\&]
\varnothing\arrow[r]\&\unsure{x|{\cdot}}\arrow[d]\\
\&y
\end{tikzcd}\qquad
\begin{tikzcd}[ampersand replacement=\&]x'_{\rm cof}\arrow[r,shift left]\arrow[r,shift right]\&y'_{\rm cof}\arrow[r,line width=1.3pt]\&y\end{tikzcd}}
Si ce lemme est valable, on trouve un foncteur $\deleted{\unsure{W}}C\longrightarrow W_{\mathrm{Cof}}^{-1}\,\mathrm{Cof}$, car la compatibilité de $u\mapsto\Phi(u)$ aux identités et aux compositions est immédiate.
% archive sheet number: 1

\scanpage{3}
\mspage{0001}

% printed form
\noindent\texttt{DATA SET UTILITY - GENERATE}\hfill\texttt{PAGE 0001}\\
\texttt{IEB352I WARNING : OUTPUT RECFM/LRECL/BLKSIZE COPIED FROM INPUT}\\
\texttt{PROCESSING ENDED AT EOD}

\handwritten{\diagram{\begin{tikzcd}[ampersand replacement=\&,column sep=small,row sep=small]
x \arrow[r,tail,"\sim"] \arrow[d,"\sim"'] \& x' \arrow[dl,"\sim"] \\
y \arrow[ur,bend right,"\sim"] \& \deleted{\unsure{g}} \\
\mathrm{cof} \arrow[u] \& {}
\end{tikzcd}}}

\handwritten{\diagram{\begin{tikzcd}[ampersand replacement=\&,column sep=small,row sep=small]
{} \& y \arrow[d,"\sim"] \arrow[dr,"\sim"] \& {} \\
x \arrow[r] \& y' \arrow[dr] \arrow[d,tail,no head] \& y'' \arrow[d,"\sim"] \\
{} \& {} \& y'''
\end{tikzcd}}}

\handwritten{\diagram{\begin{tikzcd}[ampersand replacement=\&,column sep=small,row sep=small]
a \arrow[dr,tail] \arrow[dd,tail] \& y \arrow[d,"\sim"] \\
{} \& \bar y \arrow[dl,"\sim"] \arrow[d,"\sim"] \\
\bar y' \arrow[r,tail] \& \bar y''
\end{tikzcd}}}

\handwritten{\diagram{\begin{tikzcd}[ampersand replacement=\&,column sep=small,row sep=small]
a \arrow[r,tail] \arrow[d,equal] \& \unsure{z|y} \arrow[d] \& y \arrow[l] \arrow[d,equal] \\
x \arrow[r,tail] \arrow[d,equal] \& y'' \& y \arrow[l] \arrow[d,equal] \\
x \arrow[r] \& \bar y' \arrow[u] \& y \arrow[l]
\end{tikzcd}}}

\handwritten{\diagram{\begin{tikzcd}[ampersand replacement=\&,column sep=small,row sep=small]
{} \arrow[r,tail,no head] \& {}
\end{tikzcd}}}

\handwritten{\diagram{\begin{tikzcd}[ampersand replacement=\&,column sep=small,row sep=small]
{} \arrow[r,tail] \arrow[d,"1"'] \& {} \arrow[d,"\sim"] \& {} \arrow[l,"\sim"'] \arrow[d,"1"] \\
{} \arrow[r,tail] \& {} \& {} \arrow[l] \\
{} \arrow[u,"1"'] \arrow[r,tail] \& {} \arrow[u] \& {} \arrow[l] \arrow[u,"1"]
\end{tikzcd}}}

\handwritten{\diagram{\begin{tikzcd}[ampersand replacement=\&,column sep=small,row sep=small]
x \arrow[r,tail] \& {}
\end{tikzcd}}}

\handwritten{\deleted{$x\amalg x$.}}

\handwritten{\diagram{\begin{tikzcd}[ampersand replacement=\&,column sep=small,row sep=small]
{} \& {} \& y \arrow[dl,"\sim"'] \arrow[d,"r"] \arrow[dr,"\sim"] \& {} \& {} \\
a \arrow[r,tail,"f"] \arrow[drr,"{\unsure{i_1 f|if}}"'] \& y_1 \& {} \& y_2 \& x \arrow[l,tail,"g"'] \arrow[dll,"{i_2g}"] \\
{} \& {} \& \bar y \arrow[ul,"{r_1}"] \arrow[ur,"{r_2}"'] \& {} \& {}
\end{tikzcd}}}

\handwritten{\diagram{\begin{tikzcd}[ampersand replacement=\&,column sep=small,row sep=small]
x \arrow[r] \& y
\end{tikzcd}}}

\handwritten{\diagram{\begin{tikzcd}[ampersand replacement=\&,column sep=small,row sep=small]
x \arrow[r,"{\deleted{\unsure{u}}}"] \arrow[dr,"u"] \arrow[ddr,"{u'}"'] \& y \arrow[d,equal] \\
{} \& y \arrow[d,equal] \\
{} \& {}
\end{tikzcd}}}

\handwritten{\diagram{\begin{tikzcd}[ampersand replacement=\&,column sep=small,row sep=small]
x \arrow[r,shift left,"u"] \arrow[r,shift right] \& y
\end{tikzcd}}}

\handwritten{\diagram{\begin{tikzcd}[ampersand replacement=\&,column sep=small,row sep=small]
x \arrow[r,"f"] \& y
\end{tikzcd}}}

\handwritten{\diagram{\begin{tikzcd}[ampersand replacement=\&,column sep=small,row sep=small]
x \arrow[r] \& \bar y \& y \arrow[l]
\end{tikzcd}}}

\handwritten{\diagram{\begin{tikzcd}[ampersand replacement=\&,column sep=small,row sep=small]
x \arrow[r,"g"] \& y
\end{tikzcd}}}

\handwritten{\diagram{\begin{tikzcd}[ampersand replacement=\&,column sep=small,row sep=small]
x \arrow[r,"f"] \arrow[d,equal] \& y \arrow[d,tail,"i"] \& y \arrow[l,"1"'] \arrow[d,equal] \\
x \arrow[r,"h"] \arrow[d,equal] \& \bar y \& y \arrow[l,"i"'] \arrow[d,equal] \\
x \arrow[r,"g"] \& y \arrow[u,tail,"i"] \& y \arrow[l,"1"']
\end{tikzcd}}}

\handwritten{\diagram{\begin{tikzcd}[ampersand replacement=\&,column sep=small,row sep=small]
x \arrow[r,shift left,"f"] \arrow[r,shift right,"g"'] \arrow[dr,"h"'] \& y \arrow[d,"i","\sim"'] \\
{} \& \bar y
\end{tikzcd}}}

\handwritten{\diagram{\begin{tikzcd}[ampersand replacement=\&,column sep=small,row sep=small]
{} \arrow[r] \& {}
\end{tikzcd}}}

\handwritten{\diagram{\begin{tikzcd}[ampersand replacement=\&,column sep=small,row sep=small]
x \arrow[r,"\sim"] \arrow[d] \& y \\
{} \arrow[ur] \& {}
\end{tikzcd}}}

\handwritten{\diagram{\begin{tikzcd}[ampersand replacement=\&,column sep=small,row sep=small]
y\vee y \arrow[r] \& y \& {} \\
x\vee x \arrow[u,"{f\vee g}"] \arrow[r] \& C \arrow[u,"u"'] \arrow[r,shift left,"{\deleted{\unsure{g}}}"] \arrow[r,shift right] \& \deleted{\unsure{y}}x
\end{tikzcd}}}

\handwritten{\diagram{\begin{tikzcd}[ampersand replacement=\&,column sep=small,row sep=small]
x \arrow[r,shift left,"{\unsure{h_1|k_1}}"] \arrow[r,shift right,"{g'_2}"'] \& C \arrow[r,"{\deleted{g}h}"] \& y
\end{tikzcd}}}

\handwritten{$a.\qquad x$}

\handwritten{\diagram{\begin{tikzcd}[ampersand replacement=\&,column sep=small,row sep=small]
a \arrow[dr] \& y \arrow[d,"\sim"] \arrow[dr] \& \bar{\unsure{z|x|{\cdot}}} \arrow[d] \\
{} \& \bar y \arrow[dr] \& {} \arrow[d,"\sim"] \\
{} \& {} \& {}
\end{tikzcd}}}

\handwritten{\diagram{\begin{tikzcd}[ampersand replacement=\&,column sep=small,row sep=small]
a \arrow[r] \arrow[d,tail] \& \unsure{x|a} \arrow[r,tail,"\sim"] \arrow[d,"\sim"'] \& \bar x \arrow[dl,"\sim","{\unsure{w|u}}"'] \\
b \arrow[r] \arrow[urr] \& y \& {}
\end{tikzcd}}}

\handwritten{\diagram{\begin{tikzcd}[ampersand replacement=\&,column sep=small,row sep=small]
\unsure{x|{\cdot}} \arrow[dr] \& {} \arrow[d] \\
{} \& {}
\end{tikzcd}}}

\handwritten{\diagram{\begin{tikzcd}[ampersand replacement=\&,column sep=small,row sep=small]
a \arrow[r] \& y \arrow[d,tail] \\
{} \& {}
\end{tikzcd}}}

\handwritten{$x$}

\handwritten{\diagram{\begin{tikzcd}[ampersand replacement=\&,column sep=small,row sep=small]
y \& {} \arrow[l]
\end{tikzcd}}}

\handwritten{\diagram{\begin{tikzcd}[ampersand replacement=\&,column sep=small,row sep=small]
y \arrow[r,tail] \& {}
\end{tikzcd}}}

\handwritten{\diagram{\begin{tikzcd}[ampersand replacement=\&,column sep=small,row sep=small]
\overset{\rm cof}{\unsure{x|a}} \arrow[r,tail] \arrow[dr,tail] \& \overset{\rm cof}{y'} \arrow[r,"\sim"'] \arrow[d,tail,"\sim"'] \& \overset{\rm cof}{y} \arrow[dl,dashed,bend left] \\
{} \& \bar y \arrow[ur,"\sim"] \& {}
\end{tikzcd}}}

\handwritten{\diagram{\begin{tikzcd}[ampersand replacement=\&,column sep=small,row sep=small]
{} \& x \arrow[dl,no head] \& y \\
y^{\rm Cof} \arrow[r,"\sim"'] \arrow[dr,"?"] \arrow[dd,tail] \& \bar y \arrow[ur] \arrow[dr] \& {} \\
{} \& \bar y' \arrow[r] \& y_1 \\
y_2 \& {} \& {}
\end{tikzcd}}}

\handwritten{\diagram{\begin{tikzcd}[ampersand replacement=\&,column sep=small,row sep=small]
\bar y \arrow[r,"\sim"'] \arrow[d,tail,"\sim"'] \& y' \arrow[dl,bend left] \\
y'' \arrow[ur] \& {}
\end{tikzcd}}}

\handwritten{\diagram{\begin{tikzcd}[ampersand replacement=\&,column sep=small,row sep=small]
x \arrow[r] \& y'.
\end{tikzcd}}}

\handwritten{\diagram{\begin{tikzcd}[ampersand replacement=\&,column sep=small,row sep=small]
{} \& \varnothing \arrow[r,tail] \arrow[d,no head,dashed] \& y \arrow[d,tail] \arrow[dl] \\
x\amalg x \arrow[r,tail] \& \widetilde y \arrow[r,"\sim"'] \arrow[d,tail] \& \bar y \arrow[dl] \\
{} \& \widetilde y' \arrow[ur,shift right] \& {}
\end{tikzcd}}}

\handwritten{\diagram{\begin{tikzcd}[ampersand replacement=\&,column sep=small,row sep=small]
x \arrow[dr,tail] \& {} \& y \arrow[dl,"\sim"'] \arrow[dd,tail,"\sim"'] \arrow[dr,"\sim"] \& {} \\
{} \& y' \arrow[dr,"{\unsure{g|q}}"'] \& {} \& y'' \\
{} \& {} \& y''' \arrow[ul,"\sim"] \arrow[ur,"\sim"] \& {}
\end{tikzcd}}}

\handwritten{\diagram{\begin{tikzcd}[ampersand replacement=\&,column sep=small,row sep=small]
x \& y \arrow[d,tail] \\
x \arrow[r,shift left] \arrow[r,shift right] \& \bar y
\end{tikzcd}}}

\handwritten{\diagram{\begin{tikzcd}[ampersand replacement=\&,column sep=small,row sep=small]
{} \& {} \& y \arrow[d,tail] \\
x \arrow[r,tail] \& \bar y' \arrow[r,"\sim"'] \& \bar y
\end{tikzcd}}}

\handwritten{\diagram{\begin{tikzcd}[ampersand replacement=\&,column sep=small,row sep=small]
{} \& {} \& y \arrow[d,tail,"\sim"] \\
\deleted{\unsure{x\amalg x}} \arrow[r,tail] \& \deleted{\unsure{x\amalg x}}\bar y' \arrow[r,"\sim"'] \& \bar y
\end{tikzcd}}}

% archive sheet number: 2

\scanpage{4}

\mspage{\uline{2}}

\deleted{\diagram{\begin{tikzcd}[ampersand replacement=\&,column sep=small,row sep=small]
\tilde x \arrow[r,"{\deleted{\unsure{\varphi}}}"] \arrow[rr,bend left,"\varphi"] \arrow[d,"\sim"'] \& \tilde y \arrow[r,tail,"i","\sim"'] \arrow[d,"\sim"] \& \tilde y' \arrow[rr,bend left,"\psi"] \& \tilde z \arrow[r,tail,"j","\sim"'] \arrow[d,"\sim"'] \& \tilde z' \arrow[dl] \\
x \arrow[r,"f",no head] \& y \arrow[rr,"g"] \& {} \& z \& {}
\end{tikzcd}}}

\deleted{\diagram{\begin{tikzcd}[ampersand replacement=\&,column sep=small,row sep=small]
\tilde x \arrow[dr] \& \tilde y \arrow[d,"i"] \arrow[dr,"\psi"] \& \tilde z \arrow[d,"j"] \\
{} \& \tilde y' \& \tilde z'
\end{tikzcd}}}

\diagram{\begin{tikzcd}[ampersand replacement=\&,column sep=small,row sep=small]
\tilde x \arrow[r,"{\deleted{\unsure{\tilde f|\tilde\varphi}}}"] \arrow[dr,"{\tilde f}"'] \arrow[ddd,dashed,"\sim"'] \& \tilde y \arrow[d,tail,"{i_f}"] \arrow[dr,"{\tilde g}"] \& \tilde z \arrow[d,tail,"{i_g}"] \\
{} \& \tilde y' \arrow[dd,dashed,"\sim"'] \arrow[dr,no head,"{\deleted{\unsure{u}}\tilde g'}"] \& \tilde z' \arrow[d,tail] \\
{} \& {} \& \tilde z'' \arrow[d,dashed,no head] \\
x \arrow[r,"f"] \& y \arrow[r,"g"] \& z
\end{tikzcd}}

\deleted{\unsure{On|Ce}} Montrons que ce foncteur transforme $f\in W$ en iso en effet, si $f\in W$, donc $f\alpha_x\in W$, on voit aussitôt qu'on aura $\widetilde f\in W$, donc $[\widetilde f]$ est un iso.

Considérons les deux composés

\diagram{\begin{tikzcd}[ampersand replacement=\&,column sep=small,row sep=small]
W_{\rm Cof}^{-1}\,\mathrm{Cof} \arrow[r,shift left,"\Psi"] \& W^{-1}C \arrow[l,shift left,"\Phi"]
\end{tikzcd}}

\deleted{pour voir} \inserted{prouvons} qu'\deleted{ils sont} sont \deleted{\unsure{égaux}} iso aux foncteurs identiques.

Pour $\Phi\Psi$, c'est $=$ l'identité, si on prend soin de choisir $\widetilde x\xrightarrow{\alpha_x}x$ l'identité, quand $x\in\mathrm{Cof}$. Pour $\Psi\Phi$, il \deleted{suffit de} \inserted{\unsure{veut|vaut|faut}} voir que le foncteur \deleted{$\Psi\Phi\overset{\rm déf}{=}$} $R\overset{\rm déf}{=}\Psi\Phi\mid C$

\[
 x\longmapsto\widetilde x\qquad C\longrightarrow W^{-1}C
\]

est iso \deleted{à l'identité} \inserted{au foncteur can}, on a

\diagram{\begin{tikzcd}[ampersand replacement=\&,column sep=small,row sep=small]
R(x) \arrow[d,equal] \arrow[r,"{[\alpha_x]}"] \& x \\
\widetilde x \& {}
\end{tikzcd}}

dans $C$ donc dans $W^{-1}C$ (soit $[\alpha_x]$ qui est un iso, il faut vérifier que c'est fonctoriel

\diagram{\begin{tikzcd}[ampersand replacement=\&,column sep=small,row sep=small]
\widetilde x \arrow[d,"{\alpha_x}"'] \arrow[rr,bend left,"{\widetilde f}"] \& \widetilde y \arrow[r,tail,"{i_f}"'] \arrow[d,"{\alpha_y}"] \& \widetilde y' \arrow[dl,"{\beta_f}"] \\
x \arrow[r,"f"] \& y \& {}
\end{tikzcd}}

i.e qu'on a
\[
 [\alpha_y]\deleted{\unsure{{}\circ}}\bigl([i_f]^{-1}\deleted{\unsure{{}[\,]}}[\tilde f]\bigr)=[f][\alpha_x]
\]
Mais $\alpha_y=\beta_f i_f$ d'où le premier mb est $[\beta_f\tilde f]=[f\alpha_x]=[f][\alpha_x]$ cqfd.
% archive sheet number: 3

\scanpage{5}
% printed form
{\fontsize{4}{5}\selectfont
\begin{verbatim}
//*   DOCUGRA20005        I
           00000000      00000000      00000000      00000000
         0000000000    0000000000    0000000000    0000000000
       00      0000  00      0000  00      0000  00      0000
      00     00 00  00     00 00  00     00 00  00     00 00
     00    00  00  00    00  00  00    00  00  00    00  00
    00   00   00  00   00   00  00   00   00  00   00   00
   00  00    00  00  00    00  00  00    00  00  00    00
  00 00     00  00 00     00  00 00     00  00 00     00
 0000      00  0000      00  0000      00  0000      00
000       00  000       00  000       00  000       00
0000000000    0000000000    0000000000    0000000000
00000000      00000000      00000000      00000000
\end{verbatim}
}

{\fontsize{4}{5}\selectfont
\begin{verbatim}
DDDDDDDDD     OOOOOOOOOOOO   CCCCCCCCCC   UU        UU   GGGGGGGGGG   RRRRRRRRRRR    AAAAAAAAAA    2222222222
DDDDDDDDDD    OOOOOOOOOOOO  CCCCCCCCCCCC  UU        UU  GGGGGGGGGGGG  RRRRRRRRRRRR  AAAAAAAAAAAA  222222222222
DD       DD   OO        OO  CC        CC  UU        UU  GG        GG  RR        RR  AA        AA  22        22
DD        DD  OO        OO  CC            UU        UU  GG            RR        RR  AA        AA            22
DD        DD  OO        OO  CC            UU        UU  GG            RR        RR  AA        AA            22
DD        DD  OO        OO  CC            UU        UU  GG            RRRRRRRRRRRR  AAAAAAAAAAAA           22
DD        DD  OO        OO  CC            UU        UU  GG     GGGGG  RRRRRRRRRRR   AAAAAAAAAAAA         22
DD        DD  OO        OO  CC            UU        UU  GG     GGGGG  RR    RR      AA        AA       22
DD        DD  OO        OO  CC            UU        UU  GG        GG  RR     RR     AA        AA     22
DD       DD   OO        OO  CC        CC  UU        UU  GG        GG  RR      RR    AA        AA   22
DDDDDDDDDD    OOOOOOOOOOOO  CCCCCCCCCCCC  UUUUUUUUUUUU  GGGGGGGGGGGG  RR       RR   AA        AA  222222222222
DDDDDDDDD     OOOOOOOOOOOO   CCCCCCCCCC    UUUUUUUUUU    GGGGGGGGGG   RR        RR  AA        AA  222222222222
\end{verbatim}
}

{\fontsize{4}{5}\selectfont
\begin{verbatim}
****A  START  JOB    5  DOCUGRA2                        ROOM        2.24.59 PM 24 AUG 82  R6.PR1    SYS JES2  JOB    5  START  A****
****A  START  JOB    5  DOCUGRA2                        ROOM        2.24.59 PM 24 AUG 82  R6.PR1    SYS JES2  JOB    5  START  A****
****A  START  JOB    5  DOCUGRA2                        ROOM        2.24.59 PM 24 AUG 82  R6.PR1    SYS JES2  JOB    5  START  A****
****A  START  JOB    5  DOCUGRA2                        ROOM        2.24.59 PM 24 AUG 82  R6.PR1    SYS JES2  JOB    5  START  A****
****A  START  JOB    5  DOCUGRA2                        ROOM        2.24.59 PM 24 AUG 82  R6.PR1    SYS JES2  JOB    5  START  A****
****A  START  JOB    5  DOCUGRA2                        ROOM        2.24.59 PM 24 AUG 82  R6.PR1    SYS JES2  JOB    5  START  A****
DERNIERE MISE A JOUR DE JOURNAL : 11-08-82
\end{verbatim}
}
% archive sheet number: 4

\scanpage{6}

\mspage{3-}
En résumé, on trouve

\uline{Prop.} Pour que $\Psi:W_{\rm Cof}^{-1}C_{\rm Cof}\longrightarrow W^{-1}C$ soit une équiv. de cat, il f. et s. que la condition suivante soit vérifiée :

$(\ast)$ Pour \deleted{tout} diagrammes

\diagram{\begin{tikzcd}[ampersand replacement=\&,column sep=small,row sep=small]
x \arrow[r,shift left,"\varphi"] \arrow[r,shift right,"\psi"'] \& y \arrow[r,"f","\sim"'] \& y'
\end{tikzcd}}

avec $f\varphi=f\psi$, $f\in W$, $x,y\in C_{\rm Cof}$, \inserted{(d'où $[\varphi]=[\psi]$ dans $W^{-1}C$),} on a $[\varphi]=[\psi]$ dans $W_{\rm Cof}^{-1}C_{\rm Cof}$.

(condition impliquée par celle que $\Psi$ soit \uline{fidèle})

Nous allons travailler maintenant avec $C_{\rm Cof}$ au lieu de $C$. -- ou quitte à changer de notations, on va supposer que tt objet de $C$ est cofibrant [\uline{NB} $C_{\rm Cof}$ avec les notions induites par $W$ et $\mathrm{Cof}$ est encore une catégorie \unsure{-1|nt|C1} : C3\,$]$, l'inclusion de $C_{\rm Cof}$ dans $C$ commute aux sommes amalgamées \deleted{qui} relatives : des diagrammes

\diagram{\begin{tikzcd}[ampersand replacement=\&,column sep=small,row sep=small]
a \arrow[r,"f"] \arrow[d,"g"'] \& b \\
c \& {}
\end{tikzcd}}

tels que $f$ ou $g$ soit des cof (les seules sommes dont on postule l'existence !).$]$

Soit, pour $x,y$ dans $C$,

\[
 \underline{\mathrm{Ch}}_0(x,y)=\text{catégorie des diagrammes}
\]

\diagram{\begin{tikzcd}[ampersand replacement=\&,column sep=small,row sep=small]
x \arrow[dr,"f"'] \& y \arrow[d,tail,"i"'," \sim"] \\
{} \& \bar y
\end{tikzcd}}

où $\unsure{j|i}\in W\cap\mathrm{cof}$ \deleted{\unsure{triv|catégorie}} (équiv. faibles \inserted{si $\unsure{\mathrm{cf}_t|\mathrm{cof}_t}=\mathrm{cof}_{W}$} couniverselles). On définit des accouplements

\[
 \underline{\mathrm{Ch}}_0(x,y)\times\underline{\mathrm{Ch}}_0(y,z)\longrightarrow
 \underline{\mathrm{Ch}}_0(x,z),\quad (F,G)\longmapsto G\circ F\qquad\text{par}
\]

\diagram{\begin{tikzcd}[ampersand replacement=\&,column sep=small,row sep=small]
a \arrow[dr,"f"] \& y \arrow[d,tail,"i","\sim"'] \arrow[dr,"g"] \& z \arrow[d,tail,"j","\sim"'] \\
{} \& \bar y \arrow[dr,"{g'}"'," {\text{coca}\deleted{\unsure{g}}\text{rt}}"] \& \bar z \arrow[d,tail,"k","\sim"'] \\
{} \& {} \& \bar z'
\end{tikzcd}}

$\deleted{\unsure{gf}}\quad (g,j)\circ(f,i)=(g'g)\mathord{.}kj$
% archive sheet number: 5

\scanpage{7}
% printed form
{\fontsize{4}{5}\selectfont
\begin{verbatim}
PPPPPPPPPPPP  LL            OOOOOOOOOOOO  TTTTTTTTTTTT  EEEEEEEEEEEE  DDDDDDDDDD     IIIIIIIIII   TTTTTTTTTTTT
PP        PP  LL            OO        OO       TT       EE            DD       DD        II            TT
PP        PP  LL            OO        OO       TT       EE            DD        DD       II            TT
PP        PP  LL            OO        OO       TT       EE            DD        DD       II            TT
PPPPPPPPPPPP  LL            OO        OO       TT       EEEEEEEE      DD        DD       II            TT
PPPPPPPPPPP   LL            OO        OO       TT       EEEEEEEE      DD        DD       II            TT
PP            LL            OO        OO       TT       EE            DD        DD       II            TT
PP            LL            OO        OO       TT       EE            DD        DD       II            TT
PP            LL            OO        OO       TT       EE            DD       DD        II            TT
PP            LLLLLLLLLLLL  OOOOOOOOOOOO       TT       EEEEEEEEEEEE  DDDDDDDDDD     IIIIIIIIII        TT
PP            LLLLLLLLLLLL  OOOOOOOOOOOO       TT       EEEEEEEEEEEE  DDDDDDDDD      IIIIIIIIII        TT

****A   END   JOB 7747  PLOTEDIT                        ROOM        8.39.37 AM 25 AUG 82  R6.PR1    SYS JES2  JOB 7747   END   A****
****A   END   JOB 7747  PLOTEDIT                        ROOM        8.39.37 AM 25 AUG 82  R6.PR1    SYS JES2  JOB 7747   END   A****
****A   END   JOB 7747  PLOTEDIT                        ROOM        8.39.37 AM 25 AUG 82  R6.PR1    SYS JES2  JOB 7747   END   A****
****A   END   JOB 7747  PLOTEDIT                        ROOM        8.39.37 AM 25 AUG 82  R6.PR1    SYS JES2  JOB 7747   END   A****
****A   END   JOB 7747  PLOTEDIT                        ROOM        8.39.37 AM 25 AUG 82  R6.PR1    SYS JES2  JOB 7747   END   A****
****A   END   JOB 7747  PLOTEDIT                        ROOM        8.39.37 AM 25 AUG 82  R6.PR1    SYS JES2  JOB 7747   END   A****
\end{verbatim}
}
% archive sheet number: 6

\scanpage{8}

\mspage{4)}
Cette composition est associative à iso can. près, et elle a des identités bilatères (\deleted{\unsure{pour|par}} \unsure{un tel|cet|unique|un} iso. can. !) savoir

\diagram{\begin{tikzcd}[ampersand replacement=\&,column sep=small,row sep=small]
x \arrow[dr,"1"'] \& x \arrow[d,tail,"1","\sim"'] \\
{} \& x
\end{tikzcd}}

noté $\underline{1}_x$

On a d'autre part des \deleted{\unsure{applic}} foncteurs

\[
 \operatorname{Hom}_{C}(x,y)_{\rm discret}\longrightarrow\underline{\mathrm{Ch}}_0(x,y)
\]
\deleted{compatibles avec}
\[
 f\longmapsto (f,\mathrm{id}_x)
\]

compatibles aux compositions. Ceci posé, soit

\[
 \pi_0(x,y)\overset{\rm déf}{=}\pi_0\bigl(\underline{\mathrm{Ch}}_0(x,y)\bigr),
\]

alors la composit. des $\underline{\mathrm{Ch}}_0$ définit une composition des $\pi_0(x,y)$, qui est associative (str\textsuperscript{t} !) et admet les classes des $\underline{1}_a$ comme unités bilatères. D'où une catégorie, que je note $\widetilde C^{\unsure{2|z|?}}$, et un foncteur

\[
 C\longrightarrow\widetilde C
\]

Pour l'instant, on n'a utilisé rien sur $(C,W)$, [\deleted{sauf} que \unsure{$\subset$|la} la donnée des $\mathrm{cof}_{-}\subset\mathrm{Fl}\,C$, on convient de prendre les flèches $i$ dans \deleted{\unsure{cof triv|cofibs}} \inserted{$\mathrm{couniv.}W$} -- mais \deleted{\unsure{on}} on peut prendre la construction, une cof \inserted{pas néc. $=\mathrm{cof}_{W}$,} telle que $\mathrm{cof}\cap W\subset\deleted{\unsure{W}}\mathrm{Couniv.}\,W$\,]. \deleted{Mais,} On veut prouver que ce foncteur transforme $f\in W$ en iso.

\diagram{\begin{tikzcd}[ampersand replacement=\&,column sep=small,row sep=small]
x \arrow[r,"f","\sim"'] \& y \arrow[dr,"{g?}"'] \& x \arrow[d,tail,"{i?}","\sim"'] \\
{} \& {} \& \bar x
\end{tikzcd}}

on cherche $g,i$ tels que $(gf,i)\in\mathrm{comp.}\ \mathrm{conn.}$ de $\underline{1}_x$ dans $\pi_0(\underline{\mathrm{Ch}}_0(x,x))$ et $(\bar f g,j)\in\mathrm{comp.}\ \mathrm{conn.}$ de $\underline{1}_y$

\diagram{\begin{tikzcd}[ampersand replacement=\&,column sep=small,row sep=small]
y \arrow[dr,"g"'] \& x \arrow[r,"f","\sim"'] \arrow[d,tail,"i","\sim"'] \& y \arrow[d,tail,"j"'," \sim"] \\
{} \& \bar x \arrow[r,"{\text{(cocart)}\ \sim}","{\bar f\ \sim}"'] \& \bar y
\end{tikzcd}}

% archive sheet number: 7

\scanpage{9}
% printed form
{\fontsize{4}{5}\selectfont
\begin{verbatim}
PPPPPPPPPPP   LL            OOOOOOOOOOOO  TTTTTTTTTTTT  EEEEEEEEEEEE  DDDDDDDDD      IIIIIIIIII   TTTTTTTTTTTT
PPPPPPPPPPPP  LL            OOOOOOOOOOOO  TTTTTTTTTTTT  EEEEEEEEEEEE  DDDDDDDDDD     IIIIIIIIII   TTTTTTTTTTTT
PP        PP  LL            OO        OO       TT       EE            DD       DD        II            TT
PP        PP  LL            OO        OO       TT       EE            DD        DD       II            TT
PP        PP  LL            OO        OO       TT       EE            DD        DD       II            TT
PPPPPPPPPPPP  LL            OO        OO       TT       EEEEEEEE      DD        DD       II            TT
PPPPPPPPPPP   LL            OO        OO       TT       EEEEEEEE      DD        DD       II            TT
PP            LL            OO        OO       TT       EE            DD        DD       II            TT
PP            LL            OO        OO       TT       EE            DD        DD       II            TT
PP            LL            OO        OO       TT       EE            DD       DD        II            TT
PP            LLLLLLLLLLLL  OOOOOOOOOOOO       TT       EEEEEEEEEEEE  DDDDDDDDDD     IIIIIIIIII        TT
PP            LLLLLLLLLLLL  OOOOOOOOOOOO       TT       EEEEEEEEEEEE  DDDDDDDDD      IIIIIIIIII        TT
\end{verbatim}
}

\handwritten{$\mathscr{P}$}

{\fontsize{4}{5}\selectfont
\begin{verbatim}
****A  START  JOB 7747  PLOTEDIT                        ROOM        8.39.24 AM 25 AUG 82  R6.PR1    SYS JES2  JOB 7747  START  A****
****A  START  JOB 7747  PLOTEDIT                        ROOM        8.39.24 AM 25 AUG 82  R6.PR1    SYS JES2  JOB 7747  START  A****
****A  START  JOB 7747  PLOTEDIT                        ROOM        8.39.24 AM 25 AUG 82  R6.PR1    SYS JES2  JOB 7747  START  A****
****A  START  JOB 7747  PLOTEDIT                        ROOM        8.39.24 AM 25 AUG 82  R6.PR1    SYS JES2  JOB 7747  START  A****
****A  START  JOB 7747  PLOTEDIT                        ROOM        8.39.24 AM 25 AUG 82  R6.PR1    SYS JES2  JOB 7747  START  A****
****A  START  JOB 7747  PLOTEDIT                        ROOM        8.39.24 AM 25 AUG 82  R6.PR1    SYS JES2  JOB 7747  START  A****
DERNIERE MISE A JOUR DE JOURNAL : 11-08-82
\end{verbatim}
}

\handwritten{\begin{tabular}{ll|ll}CRAY 1 & 579/1 & 849 & CRAY 1\\GRAPHI1 & \unsure{2|L} & &\end{tabular}}

{\fontsize{4}{5}\selectfont
\begin{verbatim}
TRACEUR  I   BUFFER  1020    BANDE  9 PISTES   PLUMES :  3   XMAX =    10100.0   YMAX =    30.0   STNB =   100.00000
                                                                    25/ 8/ 82     8 HEURES 37 MINUTES
JOB              BANDE IBM NUMERO                                   TRACEUR UTILISE
                 ****************************************           REMISE DE LA BANDE
                                                                    DEBUT DE TRAITEMENT
                                                                    FIN DE TRAITEMENT
                                                                    REMISE DES GRAPH.
                                                                    INCIDENTS SURVENUS

 PLOT 3C                 CALQUE                                              0           0
CNU*0579 GRAPHI1                       IC    82237 083702                    0           0
 GRAPHI1   CNU*0579 82237 M06 P02900500        7       0                     0           0 3
\end{verbatim}
}
% archive sheet number: 8

\scanpage{10}

\mspage{5)}
On utilise le weak lifting lemma pour

\margin{\uline{NB} On utilise $y\in C_{\rm Cof}$ \deleted{\unsure{en}} ici}

\diagram{\begin{tikzcd}[ampersand replacement=\&,column sep=small,row sep=small]
\varnothing \arrow[r] \arrow[d,tail] \& x \arrow[r,tail,dashed,"i","\sim"'] \arrow[d,"f","\sim"'] \& \bar x \arrow[dl,dashed,"{\deleted{\unsure{g}}\alpha}","\sim"'] \\
y \arrow[r,"1"] \& y \arrow[ur,bend right,dashed,"g"'] \& {}
\end{tikzcd}}

donc il existe $i,\deleted{\unsure{g}}\alpha$ avec $\deleted{\unsure{g}}\alpha i=f$, \deleted{\unsure{et}} $i\in W\cap\mathrm{cof}$, et $\alpha$ admet une section $g$ \inserted{$\alpha g=1_y$}.

Je dis que $(g,i)$ est l'inverse \inserted{cherché.} \deleted{\unsure{de $(f,\mathrm{id}_y)$}} \deleted{Considérons alors}

\diagram{\begin{tikzcd}[ampersand replacement=\&,column sep=small,row sep=small]
x \arrow[dd,"{1_x}"'] \arrow[dr,"{gf}"] \arrow[dddr,dashed,"{\deleted{\unsure{gf}}f}"'] \& x \arrow[d,tail,"i","\sim"'] \arrow[ddl,tail,"{1_x}","\sim"'] \\
{} \& \bar x \arrow[dd,dashed,"{\alpha}","\sim"'] \\
x \arrow[dr,dashed,"f"'] \& {} \\
{} \& y
\end{tikzcd}}

Pour $(gf,i)\sim\underline{1}_x$, on a envie d'utiliser le diagramme commutatif ci-contre, le seul ennui c'est que, on a $\alpha\in W$ mais pas néc. $\alpha\in\mathrm{cof}$, donc $\alpha i$ $(=f)$ pas néc. dans $\deleted{\unsure{W|\mathrm{cof}}}\mathrm{cof}$ (alors qu'il est dans $W$). Ça marcherait si on pouvait trouver $y\xrightarrow[\sim]{\beta}y'$ \inserted{dans $W$,} tel que le composé $\beta f$ soit non seulement dans $W$, mais de plus, dans cof.

Visiblement, ce n'est pas tjrs le cas (prendre le cas où les cof. sont les monos, et où $f$ n'est \uline{pas} un mono !). Ça se présente mal !

\diagram{\begin{tikzcd}[ampersand replacement=\&,column sep=small,row sep=small]
x \arrow[d,"f"'] \arrow[ddr,tail] \& {} \\
y \arrow[dr,"\beta"'] \& {} \\
{} \& y'
\end{tikzcd}}

On va alors plutôt définir d'abord le foncteur canonique

\[
\begin{aligned}
 \widetilde C&\longrightarrow W^{-1}C\\
 &\unsure{\varphi|\colon|{}}\ x\longmapsto x\\
 &(f,i)\longmapsto [i]^{-1}[f]
\end{aligned}
\]

il est immédiat que c'est compatible avec les identités et la composition. Donc on a bien foncteur, \deleted{surjectif} \inserted{bijectif} sur objets. Je dis que c'est surjectif sur les flèches ; Il suffit de voir, par composition, que si $f:x\longrightarrow y$ est dans $W$, alors $[f]^{-1}$ provient de $\widetilde C$. On va prendre \deleted{$(y,x)$} $y\longrightarrow x$ de $\widetilde C$, défini par $(g,i)$, \deleted{À prouver que} \inserted{i.e.} $[i]^{-1}[g]$, pour prouver que $[i]^{-1}[g]=[f]^{-1}$, il suffit de prouver que $[f][i]^{-1}[g)=[1_y]$, or $\alpha i=f$ entraîne $[\alpha][i]=[f]$ d'où $[f][i]^{-1}=[\alpha]$ donc le pr. mb est $[\alpha][g]=[\alpha g]=[1_y]$, OK.
% archive sheet number: 9

\scanpage{11}
% printed form
{\fontsize{4}{5}\selectfont
\begin{verbatim}
//*   PLOTEDIT7747        I
           00000000      00000000      00000000      00000000
         0000000000    0000000000    0000000000    0000000000
       00      0000  00      0000  00      0000  00      0000
      00     00 00  00     00 00  00     00 00  00     00 00
     00    00  00  00    00  00  00    00  00  00    00  00
    00   00   00  00   00   00  00   00   00  00   00   00
   00  00    00  00  00    00  00  00    00  00  00    00
  00 00     00  00 00     00  00 00     00  00 00     00
 0000      00  0000      00  0000      00  0000      00
000       00  000       00  000       00  000       00
0000000000    0000000000    0000000000    0000000000
00000000      00000000      00000000      00000000

PPPPPPPPPPP   LL            OOOOOOOOOOOO  TTTTTTTTTTTT  EEEEEEEEEEEE  DDDDDDDDD      IIIIIIIIII   TTTTTTTTTTTT
\end{verbatim}
}
% archive sheet number: 10

\scanpage{12}

\mspage{6)}
La difficulté, c'est qu'il n'est pas clair que le foncteur

\[
 \widetilde C\longrightarrow W^{-1}C
\]

précédent soit aussi \uline{fidèle} (ce qui impliquerait que c'est une équivalence de catégories \inserted{et $=$ un iso}). \deleted{Dans ce cas, l'iso. inverse serait forcément donné par} \inserted{Il revient au $=$ de dire} que le foncteur can $C\longrightarrow\widetilde C$ (dont le composé avec $\widetilde C\longrightarrow W^{-1}C$ est bien le foncteur can. $C\longrightarrow W^{-1}C$) transforme flèches de $W$ en iso, ou encore que la construction qu'on vient \deleted{\unsure{de}} d'esquisser, pour \deleted{\unsure{le voir|les inv.}} \inserted{trouver} un inverse de $[f,i]$ dans $\widetilde C$, est bien valable, \deleted{\unsure{ce que}} Cela entraînerait d'ailleurs \inserted{semble-t-il,} que, \deleted{le} posant $W_0=W\cap\mathrm{cof}$, le foncteur

\[
 W_0^{-1}C\longrightarrow W^{-1}C
\]

serait une équivalence de catégories, car

\[
 \widetilde C\longrightarrow W^{-1}C
\]

se factorise visiblement par

\diagram{\begin{tikzcd}[ampersand replacement=\&,column sep=small,row sep=small]
\widetilde C \arrow[r] \arrow[d,"{\unsure{\wr|{\simeq}}}"{description},draw=none] \& W_0^{-1}C \arrow[r] \& W^{-1}C \\
W^{-1}C \& {} \& {}
\end{tikzcd}}

Mais je doute qu'il suffise d'inverser les équiv. faibles \uline{cofibrantes}, pour que ttes les équiv. faibles deviennent inversibles !

Je rectifie \deleted{donc} le tir, en prenant plutôt la catégorie (notée pourtant par $\widetilde C$) ayant $\mathord{=}^{\rm s}$ objets que $C$, \deleted{et où on désigne les flèches de $x$ à $y$} \inserted{et définie via les $\pi_0$ des $\underline{\mathrm{Ch}}_0(x,y)$}, \deleted{\unsure{définis|définies}} \inserted{cette catégorie} étant maintenant (pour $x,y$ donnés) \unsure{comme|connue} celle des diagrammes \deleted{\unsure{essayer}}

\diagram{\begin{tikzcd}[ampersand replacement=\&,column sep=small,row sep=small]
x \arrow[dr,tail,"f"'] \& y \arrow[d,"i"'," \sim"] \\
{} \& \bar y
\end{tikzcd}}

avec $\left\{\begin{aligned} f&\in\mathrm{cof}\\ i&\in W\end{aligned}\right.$
% archive sheet number: 11

\scanpage{13}

% printed form
{\fontsize{4}{5}\selectfont
\begin{verbatim}
//*   DOCUGRA20005        I
           00000000      00000000      00000000      00000000
         0000000000    0000000000    0000000000    0000000000
       00      0000  00      0000  00      0000  00      0000
      00     00 00  00     00 00  00     00 00  00     00 00
     00    00  00  00    00  00  00    00  00  00    00  00
    00   00   00  00   00   00  00   00   00  00   00   00
   00  00    00  00  00    00  00  00    00  00  00    00
  00 00     00  00 00     00  00 00     00  00 00     00
 0000      00  0000      00  0000      00  0000      00
000       00  000       00  000       00  000       00
0000000000    0000000000    0000000000    0000000000
00000000      00000000      00000000      00000000

DDDDDDDDD     OOOOOOOOOOOO   CCCCCCCCCC   UU        UU   GGGGGGGGGG   RRRRRRRRRRR    AAAAAAAAAA    2222222222
DDDDDDDDDD    OOOOOOOOOOOO  CCCCCCCCCCCC  UU        UU  GGGGGGGGGGGG  RRRRRRRRRRRR  AAAAAAAAAAAA  222222222222
DD       DD   OO        OO  CC        CC  UU        UU  GG        GG  RR        RR  AA        AA  22        22
DD        DD  OO        OO  CC            UU        UU  GG            RR        RR  AA        AA            22
DD        DD  OO        OO  CC            UU        UU  GG            RR        RR  AA        AA            22
DD        DD  OO        OO  CC            UU        UU  GG            RRRRRRRRRRRR  AAAAAAAAAAAA           22
DD        DD  OO        OO  CC            UU        UU  GG     GGGGG  RRRRRRRRRRR   AAAAAAAAAAAA         22
DD        DD  OO        OO  CC            UU        UU  GG     GGGGG  RR    RR      AA        AA       22
DD        DD  OO        OO  CC            UU        UU  GG        GG  RR     RR     AA        AA     22
DD       DD   OO        OO  CC        CC  UU        UU  GG        GG  RR      RR    AA        AA   22
DDDDDDDDDD    OOOOOOOOOOOO  CCCCCCCCCCCC  UUUUUUUUUUUU  GGGGGGGGGGGG  RR       RR   AA        AA  222222222222
DDDDDDDDD     OOOOOOOOOOOO   CCCCCCCCCC    UUUUUUUUUU    GGGGGGGGGG   RR        RR  AA        AA  222222222222

****A   END   JOB    5  DOCUGRA2                        ROOM        2.28.52 PM 24 AUG 82  R6.PR1    SYS JES2  JOB    5   END   A****
****A   END   JOB    5  DOCUGRA2                        ROOM        2.28.52 PM 24 AUG 82  R6.PR1    SYS JES2  JOB    5   END   A****
****A   END   JOB    5  DOCUGRA2                        ROOM        2.28.52 PM 24 AUG 82  R6.PR1    SYS JES2  JOB    5   END   A****
****A   END   JOB    5  DOCUGRA2                        ROOM        2.28.52 PM 24 AUG 82  R6.PR1    SYS JES2  JOB    5   END   A****
****A   END   JOB    5  DOCUGRA2                        ROOM        2.28.52 PM 24 AUG 82  R6.PR1    SYS JES2  JOB    5   END   A****
****A   END   JOB    5  DOCUGRA2                        ROOM        2.28.52 PM 24 AUG 82  R6.PR1    SYS JES2  JOB    5   END   A****
\end{verbatim}
}

\handwritten{\diagram{\begin{tikzcd}[ampersand replacement=\&,column sep=small,row sep=small]
x \arrow[dr] \& x \arrow[d,tail,"\sim"'] \\
\deleted{\unsure{x}} \arrow[r,"{\deleted{\unsure{\sim}}}"] \& \bar x
\end{tikzcd}}}

\handwritten{\diagram{\begin{tikzcd}[ampersand replacement=\&,column sep=small,row sep=small]
\widetilde C \arrow[r,"{\unsure{\mathscr F|\Phi}}"] \& W^{-1}C
\end{tikzcd}}}

\handwritten{\diagram{\begin{tikzcd}[ampersand replacement=\&,column sep=small,row sep=small]
x \arrow[dr] \arrow[dd] \& x \arrow[d,tail] \arrow[ddl,bend right] \\
{} \& x' \\
x'' \arrow[ur,"{\deleted{\unsure{\sim}}}"'] \arrow[uur,bend right] \& {}
\end{tikzcd}}}

% archive sheet number: 12

\scanpage{14}

\mspage{7)}
qui se composent encore comme précédemment. On a bien \deleted{\unsure{un|sur|donc}} foncteur

\[
 \deleted{\unsure{C\longrightarrow}}\quad\widetilde C\longrightarrow W^{-1}C
\]

% slanted
\margin{Ce foncteur est bijectif sur les objets, surjectif sur les flèches, comme on vérifie aisément (il suffit pour ce dernier point de voir que pour tt $f\in W$, $f:x\longrightarrow y$, $[f]^{-1}$ provient de $\widetilde C$, \unsure{or prenant|on prend} la flèche \unsure{idem|'id'}.)}

\diagram{\begin{tikzcd}[ampersand replacement=\&,column sep=small,row sep=small]
x \arrow[r,"i"] \arrow[d,"\sim"'] \& y \\
\deleted{\unsure{y}} \arrow[ur,"{\alpha}"] \& {}
\end{tikzcd}}

\diagram{\begin{tikzcd}[ampersand replacement=\&,column sep=small,row sep=small]
y \arrow[dr,"{1_{\unsure{x|y}}}"'] \& x \arrow[d,"f","\sim"'] \\
{} \& y
\end{tikzcd}}

\deleted{\unsure{Le premier|Le processus} de ces deux foncteurs transforme bien éléments de $W$ en flèches inversibles. \deleted{On a} Pour $f:x\xrightarrow[\sim]{f}y$ dans $W$, on considère soit $F=[f]$}

\diagram{\begin{tikzcd}[ampersand replacement=\&,column sep=small,row sep=small]
\deleted{y} \arrow[dr,tail,"{\deleted{1_y}}"'] \& \deleted{x} \arrow[d,"{\deleted{f}}","{\deleted{\sim}}"'] \\
{} \& \deleted{y}
\end{tikzcd}}

\deleted{soit $G=[1_y,f]$ on trouve les composés \deleted{On a donc} $GF=[\quad{}^x$ définissons de $\mathord{=}$}

\[
 C\longrightarrow\widetilde C
\]

qui va être un peu plus délicat. Soit $f:x\longrightarrow y$ on choisit une factorisation

\diagram{\begin{tikzcd}[ampersand replacement=\&,column sep=small,row sep=small]
a \arrow[dr,tail,"{f_0}"'] \& y \arrow[dr,dotted,bend left,"i"] \& {} \\
{} \& \bar y \arrow[u,"\alpha","\sim"'] \arrow[r,tail,dotted,"{\sim\,j}"'] \& \bar y' \arrow[ul,dashed,"\beta","\sim"']
\end{tikzcd}}

donc $\bar y$ cof

et : en vertu du weak lifting lemma (et $y$ étant cofibrant) on peut compléter par des pointillés, où $\beta$ a une section $i$. On considère alors $[jf_0,i]$, il faut montrer que ça ne dépend pas des choix faits. Donc on choisit $[f',i]$

\margin{\unsure{Ou s'est ramené aux|On s'est ramené aux} factorisations $f=\deleted{\unsure{\beta}}\alpha f_0$, \inserted{avec} $f_0\in\mathrm{cof}$ et $\deleted{\unsure{\beta}}\alpha\in W$, telles que $\deleted{\unsure{\beta}}\alpha$ ait une section}

\diagram{\begin{tikzcd}[ampersand replacement=\&,column sep=small,row sep=small]
a \arrow[r,dashed,"f"] \arrow[dr,tail,"{f'}"'] \& y \arrow[d,"i","\sim"'] \\
{} \& \unsure{\widetilde y|y'} \arrow[u,bend right,dashed,"\alpha","\sim"']
\end{tikzcd}}

tel qu'il existe \deleted{\unsure{g}} une rétraction $\alpha$ de $\widetilde y$ sur $y$, telle que $\alpha f'=f$

Considérons deux \deleted{\unsure{tels|telles}} choix, il faut voir qu'ils sont dans une $=$ composante connexe de $\underline{\mathrm{Ch}}_0(x,y)$. Mais pour deux factorisations d'une flèche $f:x\longrightarrow y$ en $\alpha f_0$ et $\alpha'f'_0$, $f_0,f'_0\in\mathrm{cof}$, $\alpha,\alpha'\in W$,

\deleted{$x\longrightarrow\unsure{y}$}

On peut les coiffer par une troisième \deleted{\unsure{$(\alpha'',f)$}} $\alpha'',f''_0$ qui \deleted{provient de} \inserted{recolle} les deux comme on voit sur le diagramme

\diagram{\begin{tikzcd}[ampersand replacement=\&,column sep=small,row sep=small]
\deleted{a} \arrow[r,"\deleted{\unsure{f}}"] \arrow[dr] \& \deleted{y} \arrow[d] \\
{} \& \deleted{y''}
\end{tikzcd}}

% archive sheet number: 13

\scanpage{15}

% printed form
{\ttfamily //* \quad GRAPHI2 0421 \quad - CRAYI}

{\fontsize{4}{5}\selectfont
\begin{verbatim}
 GGGGGGGGGG   RRRRRRRRRRR    AAAAAAAAAA   PPPPPPPPPPP   HH        HH  IIIIIIIIII   2222222222
GGGGGGGGGGGG  RRRRRRRRRRRR  AAAAAAAAAAAA  PPPPPPPPPPPP  HH        HH  IIIIIIIIII  222222222222
GG        GG  RR        RR  AA        AA  PP        PP  HH        HH      II      22        22
GG            RR        RR  AA        AA  PP        PP  HH        HH      II                22
GG            RR        RR  AA        AA  PP        PP  HH        HH      II                22
GG            RRRRRRRRRRRR  AAAAAAAAAAAA  PPPPPPPPPPPP  HHHHHHHHHHHH      II               22
GG     GGGGG  RRRRRRRRRRR   AAAAAAAAAAAA  PPPPPPPPPPP   HHHHHHHHHHHH      II             22
GG     GGGGG  RR    RR      AA        AA  PP            HH        HH      II           22
GG        GG  RR     RR     AA        AA  PP            HH        HH      II         22
GG        GG  RR      RR    AA        AA  PP            HH        HH      II       22
GGGGGGGGGGGG  RR       RR   AA        AA  PP            HH        HH  IIIIIIIIII  222222222222
 GGGGGGGGGG   RR        RR  AA        AA  PP            HH        HH  IIIIIIIIII  222222222222
\end{verbatim}
}

{\fontsize{4}{5}\selectfont
\begin{verbatim}
****A   END   JOB  421  GRAPHI2     - CRAY OUTPUT -     ROOM        3.04.20 PM 25 AUG 82  R6.PR1    SYS JES2  JOB  421   END   A****
****A   END   JOB  421  GRAPHI2     - CRAY OUTPUT -     ROOM        3.04.20 PM 25 AUG 82  R6.PR1    SYS JES2  JOB  421   END   A****
****A   END   JOB  421  GRAPHI2     - CRAY OUTPUT -     ROOM        3.04.20 PM 25 AUG 82  R6.PR1    SYS JES2  JOB  421   END   A****
****A   END   JOB  421  GRAPHI2     - CRAY OUTPUT -     ROOM        3.04.20 PM 25 AUG 82  R6.PR1    SYS JES2  JOB  421   END   A****
****A   END   JOB  421  GRAPHI2     - CRAY OUTPUT -     ROOM        3.04.20 PM 25 AUG 82  R6.PR1    SYS JES2  JOB  421   END   A****
****A   END   JOB  421  GRAPHI2     - CRAY OUTPUT -     ROOM        3.04.20 PM 25 AUG 82  R6.PR1    SYS JES2  JOB  421   END   A****
\end{verbatim}
}

\handwritten{\diagram{\begin{tikzcd}[ampersand replacement=\&,column sep=small,row sep=small]
x \arrow[dr,tail] \& y \\
{} \& y' \arrow[u] \& y'' \arrow[ul] \\
x\amalg x \arrow[r,tail] \arrow[ddr] \& y'\amalg y'' \arrow[u] \arrow[ddr] \& \deleted{\unsure{y|x}} \arrow[l,leftrightarrow] \\
{} \& {} \& {} \\
{} \& x \arrow[r,tail] \& \unsure{p|\rho|\bullet} \arrow[r,tail] \& y
\end{tikzcd}}}

\handwritten{\diagram{\begin{tikzcd}[ampersand replacement=\&,column sep=small,row sep=small]
{} \& x \arrow[r,tail] \arrow[d,"{\unsure{i_1|i_4}}"] \arrow[dr,phantom,"{\mathrm{cocart}}"{description}] \& y' \arrow[d] \\
\deleted{\unsure{y|g}}\,x \arrow[r,"{i_2}"] \arrow[d,tail] \& x\amalg x \arrow[r,tail] \arrow[d,tail] \& \deleted{x}y'\amalg x \arrow[d,tail] \\
y'' \arrow[r] \& x\amalg y'' \arrow[r,tail] \& y'\amalg y''
\end{tikzcd}}}



% archive sheet number: 14

\scanpage{16}

\mspage{8)}

\diagram{\begin{tikzcd}[ampersand replacement=\&,column sep=small,row sep=small]
x\vee x \arrow[r,tail,"{\deleted{\unsure{x|f|(}} f_0\vee f'_0\deleted{\unsure{g|)}}}"] \arrow[d,"{\nabla_x}"'] \arrow[dr,phantom,"{\mathrm{cocart}}"{description}] \& \bar y\amalg\bar y' \arrow[d,dashed] \arrow[drr,"{(\alpha,\alpha')}"] \& {} \& {} \\
x \arrow[r,tail,dashed,"g"] \arrow[rrr,bend right=17,"f"'] \& \widetilde y \arrow[r,dashed,"{\varphi_0}"] \arrow[rr,dashed,"{\deleted{\unsure{\varphi|\alpha}}\varphi}"'] \& \bar y'' \arrow[r,dashed,"{\alpha''}","{\sim}"'] \& y
\end{tikzcd}}


on factorise $\varphi$ en $\deleted{\unsure{\varphi|\alpha}}\alpha''\varphi_0$ avec $\varphi_0\in\mathrm{cof}$ $\deleted{\unsure{\varphi|\alpha}}\alpha''\in W$, et on prend
\[
 f_0''=\varphi_0 g
\]

Mais si $\alpha,\alpha'$ ont des sections (il suffirait que $\alpha$ ou $\alpha'$ les \unsure{aient|ait}) \deleted{\unsure{il reste donc|On montre donc|il reste une}} il est trivial que \deleted{\unsure{les autres le|il en est de|les autres la}} $\alpha''$ (quel lui \unsure{qu'il soit|aussi soit}) en a aussi une. Donc on est ramené au cas où on a une flèche $u:(f_0,\alpha)\longrightarrow(\unsure{f'_0|f_0|f''_0},\alpha')$ dans la catégorie des factorisations de $f$ en $\alpha f_0$. Donc \uline{Si} on choisit la section $\unsure{i^*|i|i'}$ pour $\alpha^*$ comme l'image par $u$ de la section $i'$ de $\alpha'$ \deleted{\unsure{par u|on}}, on gagne.

\diagram{\begin{tikzcd}[ampersand replacement=\&,column sep=small,row sep=small]
a \arrow[rr,"f"] \arrow[dr,"{f'_0}"'] \arrow[drr,"{f^*_0}"] \& {} \& y \\
{} \& \bar y' \arrow[ur,"{\alpha'}"] \arrow[r,leftrightarrow,"u"'] \& \bar y^* \arrow[u,"{\alpha^*}"]
\end{tikzcd}}

Donc il reste à prouver que \deleted{le choix de} pour une factorisation choisie de $f$ en $\alpha f_0$ comme dessus, $[f_0,i]$ ne dépend pas du choix de la section $i$. Comme $[f_0,i]=[\mathrm{id}_{\bar y},i]\circ(f,\mathrm{id}_{\bar y}]$, On est ramené à ceci :

\diagram{\begin{tikzcd}[ampersand replacement=\&,column sep=small,row sep=small]
x \arrow[rr,"{\unsure{f^{\sim}|f^{\nearrow}}}"] \arrow[dr,"{f_0}"'] \& {} \& y \arrow[dl,bend left=18,"i"] \arrow[dl,bend left=40,"{i'}"] \\
{} \& \bar y \arrow[ur,"{\alpha}"',"{\sim}"] \& {}
\end{tikzcd}}

$(\unsure{\ast_2|\ast\ast|\ast})$ Soit $x\xrightarrow[\sim]{f}y$ de $W$, tel que $f$ ait une \deleted{\unsure{(deux $i,i'$ de $W$)|(deux $i,i'(\ldots)$)|(deux $i,i'\in W$)}} section $i$, alors \inserted{la classe dans}\deleted{\unsure{pour deux telles sections $i,i'$|pour deux telles sections $i,i'$)}} \deleted{\unsure{les comp. connexes de|les composantes connexes de}} $\pi_0(\underline{\mathrm{Ch}}_0(x,y))$ \deleted{de} de $(1_x,i)$

\diagram{\begin{tikzcd}[ampersand replacement=\&,column sep=small,row sep=small]
\deleted{\unsure{x|y}}\quad x \arrow[dr,"{1_x}"'] \& y \arrow[d,"i"] \\
{} \& x
\end{tikzcd}}

ne dépend pas du choix de la section $i$.

Cette condition est (néc. et) suffis$^{\mathrm{te}}$ pour que $\widetilde C\longrightarrow W^{-1}C$ soit fidèle (donc un iso), comme on doit pouvoir vérifier aisément, en construisant bel et bien le foncteur cherché $C\longrightarrow\widetilde C$.

% archive sheet number: 15

\scanpage{17}
\mspage{1}

% printed form
{\fontsize{5}{6}\selectfont
\begin{verbatim}
PAGE 1            ON=CELMPQRSTV             08/25/82-15:01:26      CFT 1.09(06/28/82)  PAGE 1
\end{verbatim}
}

{\footnotesize
\begin{verbatim}
 1      1.       PROGRAM GRAPHI
 2      2.       DIMENSION X(10),Y(10)
 3      3.       DATA X/0.,0.,4.,9.,1.,2.,0.,0.,1.,0./,
 4               S    Y/7.,2.,0.,1.,0.,2.,0.,3.,0.,4./
 5      4.       NOL=10
 6      5.       XCM=10.
 7      6.       YCM=10.
 8      7.       N=10
 9      8.       WRITE(6,1)
10      9.       CALL TBRAF(NOL,XCM,YCM,0,0,0,0,0,0,0,X,Y,N,0,0.)
11     10.       WRITE(6,2)
12     11.     1 FORMAT(' TBRAF VA ETRE APPELE')
13     12.     2 FORMAT(' TBRAF A ETE APPELE')
14     13.       STOP
15     14.       END
\end{verbatim}
}

% archive sheet number: 16
